% Generated draft: requires mathematical and source review.
\documentclass[11pt]{article}
\usepackage[margin=25mm]{geometry}
\usepackage{fontspec}
\setmainfont{Latin Modern Roman}
\usepackage{amsmath,amssymb,mathtools,unicode-math}
\setmathfont{Latin Modern Math}
\usepackage{enumitem,needspace,xurl,hyperref}
\setlength{\parindent}{0pt}
\setlength{\parskip}{5pt}
\setlength{\emergencystretch}{3em}
\newcommand{\N}{\mathbb{N}}
\newcommand{\Z}{\mathbb{Z}}
\newcommand{\Q}{\mathbb{Q}}
\newcommand{\R}{\mathbb{R}}
\newcommand{\C}{\mathbb{C}}
\newcommand{\F}{\mathbb{F}}
\newcommand{\A}{\mathbb{A}}
\newcommand{\PP}{\mathbb{P}}
\newcommand{\E}{\mathbb{E}}
\newcommand{\eps}{\varepsilon}
\newcommand{\dd}{\,\mathrm{d}}
\begin{document}
\section*{UnsolvedMath ID: 30005818}
\textbf{Source problem code:} catalogue entry
\subsection{Definitions and mathematical statement}
For this draft, a geometric triangle means the filled, nondegenerate triangle
\[
\Delta=\operatorname{conv}\{a,b,c\}\subset\R^3,
\]
where \(a,b,c\) are distinct noncollinear points; write \(V(\Delta)=\{a,b,c\}\). This convention makes the intended formulation precise, but the supplied record does not explicitly distinguish filled triangles from their boundaries.

For each integer \(n\geq 1\), define
\[
t(n)=\max\left\{|\mathcal T|:\begin{array}{l}
X\subset\R^3,\quad |X|=n,\\
\mathcal T\text{ is a family of geometric triangles with }V(\Delta)\subseteq X,\\
\text{for all distinct }\Delta,\Delta'\in\mathcal T,\text{ either }\Delta\cap\Delta'=\varnothing,\\
\text{or }\Delta\cap\Delta'=\{v\}\text{ for some }v\in V(\Delta)\cap V(\Delta')
\end{array}\right\}.
\]
Thus intersections along an edge, crossings, and contacts at points that are not vertices of both triangles are excluded. No general-position assumption on \(X\) is stated, and the record does not require every point of \(X\) to be used.

\textbf{Question.} Is
\[
t(n)=o(n^2)\qquad(n\to\infty)?
\]
Explicitly, is it true that
\[
\forall\varepsilon>0\ \exists N\in\N\ \forall n\in\N,
\qquad n\geq N\Longrightarrow t(n)\leq\varepsilon n^2?
\]
The supplied reconstructed statement is the basis for this question. The supplied original extraction contains only its concluding question followed by unrelated problems, so that extraction alone does not establish the definition of \(t(n)\) or the triangle convention.
\subsection{Short English statement}
Among all sets of \(n\) points in three-dimensional Euclidean space, consider the largest possible family of triangles whose pairwise intersections are empty or consist of one shared vertex. Must the number of triangles divided by \(n^2\) tend to zero? This draft interprets triangles as filled and nondegenerate; that convention requires source confirmation.
\subsection{Sources}
\begin{itemize}
\item Supplied local UnsolvedMath registry record, permanent ID \(30005818\), title: Almost-Disjoint Geometric Triangles in Three-Space. Exact source problem code: \texttt{catalogue entry}. Collection: Oberwolfach Reports Open Problems; supplied proposed year: 2024. No proposer is identified for this record.
\item \emph{Discrete Geometry} (2024), Oberwolfach Reports, DOI 10.4171/owr/2024/3, pp. 137--202. Supplied report provenance: \url{https://doi.org/10.4171/owr/2024/3}. The background describes Oberwolfach Reports as documenting mathematical workshops and their open-problem discussions.
\item Karolyi and Solymosi, \emph{Almost Disjoint Triangles in 3-Space}: \url{https://doi.org/10.1007/s00454-002-2888-z}.
\item \emph{Oberwolfach Report 3/2024}: \url{https://ems.press/content/serial-article-files/48651}.
\item Literature assessment dated exactly \texttt{2026-08-22}: the supplied curation labels the problem open, says the best cited construction gives only order \(n^{3/2}\) almost-disjoint triangles, and says authoritative literature still records the \(o(n^2)\) question as unresolved. This is supplied dated curation derived from the OWR material, not an independently verified literature assessment; time-sensitive status and cited claims require independent checking.
\item The registry labels the statement \texttt{corrected\_verified} and reports that it was reconstructed and checked against the official Oberwolfach report. This verification claim is supplied metadata and has not been independently verified here.
\item The original extraction also contains an unrelated reference: Andrey Kupavskii and Dmitriy Zakharov, \emph{The right acute angles problem?}, European Journal of Combinatorics 89 (2020), 7 pp. It accompanies the extracted angle problem and is not supplied as support for the triangle question.
\end{itemize}
\end{document}
